\section{The Radio Access}\label{sec:radio}

The common idea of radio work since 2023 is to change the unit of service from the packet to the video
frame. A frame that misses its deadline is worthless even if most of its packets arrive. \Cref{tab:families} classifies the
resulting families together with those of \cref{sec:transport}.

\begin{table}[t]
\centering\footnotesize
\caption{Families of network and delivery mechanisms for XR since 2023, what they were evaluated on, and
whether users confirmed the gain. ``Proxy'' means a network or image metric stands in for QoE.}
\label{tab:families}
\setlength{\tabcolsep}{3.5pt}
\begin{tabularx}{\linewidth}{@{}>{\raggedright\arraybackslash}p{3.2cm}>{\raggedright\arraybackslash}X>{\raggedright\arraybackslash}p{3.0cm}>{\raggedright\arraybackslash}p{2.1cm}@{}}
\toprule
Family & Idea & Evaluated on & QoE validated?\\
\midrule
Frame-aware scheduling & Schedule PDU sets against a frame deadline & System-level simulation & No; PDB met\\
Coexistence and power & Limit interference; sleep between frames & System-level simulation & No\\
Network-assisted rate & Network recommends a rate & Simulation, testbed & No; outage\\
Wi-Fi~7 MLO, mmWave & Several links; beams that follow the head & Simulation, lab hardware & No; delay, loss\\
Early congestion signals & L4S/ECN, in-network dropping & Emulation, testbeds & Proxy; frame loss\\
QUIC and MoQ delivery & Streams, delivery timeouts & Emulation & Proxy; stalls\\
Selective protection & FEC on what matters, neural repair & Loss traces & Proxy; freezes\\
Perception-driven sending & Foveation, QoE-model encoding & Prototypes with users & Yes, small labs\\
Volumetric, splat delivery & Layered, progressive, shared content & Datasets, prototypes & Mostly proxy\\
\bottomrule
\end{tabularx}
\end{table}

\subsection{Making the radio frame-aware}
\emph{PDU sets} are the standard vehicle. 3GPP Release 18 defines a PDU set as the smallest group of
packets that carries one frame, together with its relative importance, and lets the radio network apply
QoS handling to it \cite{W4410582180}. Release 19 adds burst signalling, expedited transfer of
latency-critical content and guidance on importance values \cite{W7127418717}, plus radio features such
as latency-aware uplink prioritisation, delay status reports and uplink rate control through application
adaptation \cite{W7164211432}. \emph{Frame-aware schedulers} exploit this information. A PDU-set scheduler
raised XR capacity by up to 45\% at unchanged broadband throughput \cite{W4402156740}, and an uplink
least-slack-time scheduler gained 15\% capacity and cut XR packet delay by up to 33\% in a nine-cell dense
urban case \cite{W4406264804}.

\subsection{Coexistence and energy}
XR rarely has a cell to itself. Adding a 45~Mbps XR flow with a 10~ms packet delay budget (PDB) cost about
45~Mbps of broadband capacity \cite{W4385801627}, but interference from broadband users in other cells
removed about 80\% of XR capacity; restricting broadband throughput by about half restored it at a 20~ms
PDB \cite{W4387546593}. Power saving competes with delay: standardised and proposed sleep schemes saved
6--23\% of device energy without delaying traffic unacceptably \cite{W4406265335}.

\subsection{Telling the application what the network can carry}
A second family closes the loop from network to application (\cref{fig:loop}). NetXRate combines
core-network analytics with an O-RAN xApp to recommend rates, and cut outage by up to 96\% against
over-the-top adaptation in simulation \cite{W7161282618}. On a 4G O-RAN testbed, an xApp held VR latency
near its target and freed 15.8\% of resources for other users \cite{W4392175446}.

\subsection{Wi-Fi~7 and millimetre wave}
\emph{Multi-link operation} (MLO) in Wi-Fi~7 lets one device use several channels. In simulation, it
supported seven AR stations where a single link supported six, while policies that need simultaneous
access to all links suffered under contention \cite{W4376480829}; for the same total bandwidth, more and
narrower links gave lower delay than fewer wide ones \cite{W4412836976}. \emph{Millimetre wave} offers the
rates that uncompressed Quest~3 video would need (15.75--26.25~Gbps), yet tested hardware delivered
1.85--1.98~Gbps, and moderate head rotation at 45 degrees per second drove loss to 4.72\% in one
configuration \cite{W4401507876}. Dual-beam reception across coordinated access points improved reliability by up to 22.8\% on real head-motion traces \cite{W4409076649}.

These results share their assumptions: system-level simulation with 3GPP-style traffic, which
\cref{sec:measure} found too light, and a user counted as satisfied when enough packets meet the PDB. No
study has yet shown that frame-aware scheduling raises the QoE that \cref{sec:qoe} measured.
